% 第三部分 B：引言与相关工作（由 main_zh.tex 引入）。

\section{引言}\label{sec:intro}

RLVR 假设每个采样出的解答都能被检验。当检验代价高昂时，例如需要人工评分、长时间的测试套件或缓慢的仿真器，一次训练可能只负担得起部分样本的标签，其余的可以由评判模型补上。然而，一个昨天在样本上判断正确的评判模型，今天可能就判断错了，因为它所评判的策略在不断变化。

本文用 GRPO 在 MATH 上研究这一设定，沿用第二部分的设置，其中验证器只给每组 8 个解答中的 2 个打标签。我们比较利用另外 6 个解答的四种方式：丢掉它们；用冻结的 \kev{} 的判断作为奖励；用在线重新拟合输出头（基于迄今收集的标签）之后的 \kev{} 作为奖励；或者用一个形状相同、在同样标签上从头训练的输出头作为奖励。由于验证器实际上是可用的，我们还能逐步测量每种补全奖励在它所打分的解答上有多准确。

\section{相关工作}\label{sec:related}

半监督奖励学习和主动标注可以减少 RLHF 中人工标签的数量~\cite{ouyang2022training}。训练一次后就冻结的奖励模型，会随着策略漂移而被利用~\cite{gao2023scaling}。在训练中用新标签重新拟合奖励模型是 RLHF 中常用的补救办法，但要消耗标签。在新数据上重新校准分类器的输出层是应对分布偏移的标准做法~\cite{guo2017calibration}，这里把它在线地用在评判模型的指针头上。GRPO 和我们的更新方式沿用第二部分~\cite{shao2024deepseekmath,liu2026tricks}。
